\chapter{Mosquito bites}
\index{Mosquito bites}

\section*{Definition and Overview}
Mosquito bites are localised, inflammatory skin reactions that occur when a female mosquito punctures the skin of a host to feed on blood, which is required to nourish her developing eggs. During the biting process, the mosquito injects saliva containing a mixture of proteins that act as anticoagulants, vasodilators, and anaesthetics to facilitate feeding. The host's immune system reacts to these salivary antigens, triggering a localised hypersensitivity reaction characterised by itching, redness, and swelling. While the majority of mosquito bites represent a minor, self-limiting cutaneous irritation, mosquitoes serve as vectors for numerous global pathogens—including viruses (such as Dengue, Zika, Japanese Encephalitis, and Ross River viruses) and parasites (such as \textit{Plasmodium} species causing malaria)—making them the deadliest animal vector group to human health.

\section*{Epidemiology}
Mosquitoes are ubiquitous, with over 3,500 species found across almost every continent. While simple bites occur globally, the epidemiology of mosquito-borne diseases is highly regional and climate-dependent. Globally, mosquito-borne illnesses cause over 700,000 deaths annually, with malaria in sub-Saharan Africa accounting for the majority of these fatalities. In many countries, the most common mosquito-transmitted pathogen is the Ross River virus, which causes approximately 5,000 cases of epidemic polyarthritis annually, followed by the Barmah Forest virus. Recently, Japanese Encephalitis virus (JEV) has emerged in southeastern some countries, posing a significant public health risk. The risk of mosquito bites and disease transmission peaks during warm, wet summer months and in regions near wetlands or standing water.

\section*{Aetiology and Pathophysiology}
The cutaneous reaction to a mosquito bite is a complex immunological response to foreign salivary proteins injected during capillary cannulation.
\begin{itemize}
    \item \textbf{Vector feeding mechanism:} Female mosquitoes of the genera \textit{Aedes}, \textit{Anopheles}, and \textit{Culex} use specialized mouthparts (proboscis) to pierce the epidermis and seek a dermal capillary, injecting saliva containing immunogenic proteins directly into the skin.
    \item \textbf{IgE-mediated hypersensitivity:} The host's initial reaction is a Type I immediate hypersensitivity response, where salivary antigens bind to pre-existing IgE antibodies on dermal mast cells, triggering degranulation and the release of histamine.
    \item \textbf{Histamine-driven vasodilation:} Histamine release causes rapid capillary dilation, increased vascular permeability, and stimulation of local pruriceptive nerve fibres, resulting in the immediate formation of an itchy wheal (hive).
    \item \textbf{Cell-mediated delayed reaction:} A Type IV delayed-type hypersensitivity reaction follows, mediated by T-cells and IgG antibodies, which peaks 24 to 48 hours later and manifests as a firm, erythematous, pruritic papule.
    \item \textbf{Pathogen inoculation:} If the mosquito is infected, viruses (e.g. flaviviruses, alphaviruses) or protozoa are co-injected with the saliva, replicating locally before spreading systemically via the lymphatics and bloodstream.
\end{itemize}

\section*{Clinical Features}
The clinical presentation of mosquito bites ranges from typical localised skin reactions to severe allergic reactions and systemic symptoms of transmitted infections.
\begin{itemize}
    \item \textbf{Typical local reaction:} An immediate soft, white-to-red wheal peaking within 20 minutes, followed by a delayed-onset, firm, intensely itchy, red papule that peaks in 24--48 hours and slowly fades over a week.
    \item \textbf{Skeeter syndrome (severe local allergy):} An exaggerated, large-scale inflammatory reaction characterized by significant swelling (spanning up to 10 cm or more), warmth, erythema, and low-grade fever, which is most common in young children and non-immune travellers.
    \item \textbf{Excoriations and secondary infection:} Scratched bite marks that become broken, weeping, crusted, or painful, indicating secondary bacterial inoculations.
    \item \textbf{Ross River and Barmah Forest virus features:} Moderate fever, maculopapular rash, headache, and severe, symmetrical polyarthritis and joint stiffness, particularly affecting the wrists, fingers, knees, and ankles.
    \item \textbf{Encephalitic viral symptoms:} High fever, severe headache, neck stiffness, photophobia, confusion, and altered mental state indicating central nervous system infection.
\end{itemize}

\section*{Differential Diagnosis}
Cutaneous lesions and systemic symptoms following mosquito exposure must be differentiated from other bites, skin conditions, and febrile illnesses.
\begin{itemize}
    \item \textbf{Mosquito bites vs.\ other arthropod bites:} Flea bites are typically clustered on the lower legs and present with central haemorrhagic puncta. Bedbug bites often present in linear patterns of three ("breakfast, lunch, and dinner") on exposed areas. Spider bites may feature twin puncture marks and progress to necrosis.
    \item \textbf{Skeeter syndrome vs.\ cellulitis:} Skeeter syndrome develops rapidly (within hours of a bite), lacks purulence, and is primarily mediated by allergy, whereas cellulitis is a bacterial infection of the deep dermis that develops slowly over days, featuring expanding, painful erythema, heat, and systemic signs.
    \item \textbf{Viral exanthems vs.\ drug eruptions:} Differentiating the rash of Dengue or Ross River virus from drug-induced hypersensitivity rashes or other viral illnesses (e.g. influenza, measles).
    \item \textbf{Malaria vs.\ viral haemorrhagic fevers:} Distinguishing the cyclical high fevers, rigors, and splenomegaly of malaria from other tropical fevers.
\end{itemize}

\section*{When to Seek Medical Care}
Individuals should consult a general practitioner if local bite swelling increases progressively over several days, if the bite site displays signs of bacterial infection (e.g. increasing pain, spreading redness, pus), or if they develop systemic symptoms like fever, joint pain, or rash within three weeks of being bitten.

\textbf{Contact your local emergency services for:}
\begin{itemize}
    \item Signs of a severe systemic allergic reaction (anaphylaxis), including swelling of the face, lips, or throat, difficulty breathing, wheezing, or severe dizziness.
    \item Sudden high fever accompanied by severe headache, neck stiffness, confusion, seizures, or loss of consciousness.
    \item Severe lethargy, persistent vomiting, and spontaneous bleeding (e.g. nosebleeds, blood in urine or stools) following travel to dengue-endemic areas.
\end{itemize}

\section*{Diagnosis and Investigations}
A typical mosquito bite is diagnosed clinically based on history and physical appearance, but suspected mosquito-borne diseases require specific diagnostic testing.
\begin{itemize}
    \item \textbf{Clinical evaluation:} Assessing the number, distribution, and appearance of bites, along with a detailed history of travel to endemic areas and timing of onset.
    \item \textbf{Serology (ELISA):} Measuring IgM and IgG antibodies to detect recent infection with Ross River, Barmah Forest, Dengue, or Japanese Encephalitis viruses.
    \item \textbf{Polymerase Chain Reaction (PCR):} Performed on blood or urine during the early acute phase of infection to detect viral RNA before antibodies have developed.
    \item \textbf{Thick and thin blood films:} The diagnostic standard for malaria, allowing identification and quantification of \textit{Plasmodium} parasites under a microscope.
    \item \textbf{Full Blood Count (FBC):} Monitoring for leucopenia (low white blood cells) and thrombocytopenia (low platelets), which are characteristic of Dengue and other viral infections.
\end{itemize}

\section*{Management}
Management is directed at relieving pruritus, preventing secondary infection, and treating any transmitted pathogens.
\begin{itemize}
    \item \textbf{Symptomatic local care:} Washing the bite area with soap and water to remove residual salivary proteins, followed by the application of a cold compress or ice pack to vasoconstrict local vessels and reduce swelling.
    \item \textbf{Topical antipruritics:} Applying calamine lotion, 1\% hydrocortisone cream, or crotamiton cream to the bite site to suppress histamine-mediated itching and inflammation.
    \item \textbf{Oral antihistamines:} Utilizing second-generation, non-sedating antihistamines (e.g. cetirizine, loratadine) for daytime itching, or first-generation sedating antihistamines (e.g. promethazine) at night to aid sleep.
    \item \textbf{Oral antibiotics:} Prescribing a course of cephalexin or flucloxacillin if a bite site develops secondary bacterial cellulitis or impetigo due to scratching.
    \item \textbf{Specific pathogen management:} Providing targeted therapies for confirmed infections, such as artemisinin-based combination therapies (ACTs) for malaria, while viral infections are managed with supportive care (hydration, paracetamol; avoiding NSAIDs in suspected Dengue due to bleeding risk).
\end{itemize}

\section*{Complications}
The primary complications associated with mosquito bites arise from mechanical skin breakdown or vector-transmitted pathogens.
\begin{itemize}
    \item \textbf{Secondary bacterial infection:} Cellulitis, impetigo, or abscess formation resulting from introduction of skin commensal bacteria (e.g. \textit{Staphylococcus aureus}) into scratched bites.
    \item \textbf{Skeeter syndrome complications:} Significant swelling that causes compartment-like pressure, joint restriction, or cosmetic distress.
    \item \textbf{Chronic post-viral arthropathy:} Persistent, debilitating joint pain and fatigue lasting for months or years following Ross River or Barmah Forest virus infection.
    \item \textbf{Severe neurological sequelae:} Permanent cognitive deficits, paralysis, parkinsonian movement disorders, or death resulting from Japanese Encephalitis or Murray Valley Encephalitis.
    \item \textbf{Severe dengue and shock:} Dengue haemorrhagic fever and dengue shock syndrome, characterized by capillary leak, severe bleeding, and organ impairment.
\end{itemize}

\section*{Prognosis and Outcomes}
The prognosis for simple mosquito bites is excellent, with itching and swelling resolving spontaneously within 3 to 7 days without scarring. Skeeter syndrome also carries a highly favourable prognosis, resolving completely with topical steroids. The prognosis for mosquito-borne illnesses varies significantly; Ross River virus infection is self-limiting but can cause prolonged arthralgia, whereas encephalitis viruses carry a high mortality rate (up to 30\% for JEV) and a high rate of permanent neurological disability among survivors.

\section*{Prevention and Self-Care}
Preventing mosquito bites is the most effective way to avoid discomfort and reduce the risk of mosquito-borne diseases.
\begin{itemize}
    \item \textbf{Apply effective insect repellents:} Using personal repellents containing diethyltoluamide (DEET), picaridin, or oil of lemon eucalyptus (OLE) on all exposed skin surfaces, reapplying as directed.
    \item \textbf{Wear protective clothing:} Wearing loose-fitting, light-coloured, long-sleeved shirts and long trousers when outdoors. Heavy fabrics are preferable, as mosquitoes can bite through tight materials.
    \item \textbf{Eliminate breeding sites:} Emptying, scrubbing, or covering containers that hold standing water (e.g. flower pots, bird baths, tyres) around the home weekly to disrupt the mosquito breeding cycle.
    \item \textbf{Avoid peak activity hours:} Restricting outdoor activities during dawn and dusk, when most mosquito vector species are most active.
    \item \textbf{Install physical barriers:} Ensuring windows and doors are fitted with intact flyscreens, and utilizing insecticide-treated bed nets if sleeping in non-screened areas or outdoors.
\end{itemize}

\section*{Sources and Further Reading}
The following organisations provide reliable public health guidelines and information on mosquito bites and vector-borne diseases:
\begin{itemize}
    \item Australian Government Department of Health and Aged Care. \textit{Fight the Bite: Mosquito-borne disease prevention.} https://www.health.gov.au/
    \item Better Health Channel. \textit{Mosquitoes: Protecting yourself from bites.} https://www.betterhealth.vic.gov.au/
    \item Mayo Clinic. \textit{Mosquito bites: Symptoms, causes and home care.} https://www.mayoclinic.org/
    \item World Health Organization (WHO). \textit{Vector-borne diseases: Mosquitoes and global health impact.} https://www.who.int/
    \item Centers for Disease Control and Prevention (CDC). \textit{Preventing mosquito bites and travel safety guidelines.} https://www.cdc.gov/
\end{itemize}
